\section{\I{Post-processing output using \R} \label{sec:PostProcessing}}\index{Post processing}\index{Post-processing section}

\R\ (\url{https://www.r-project.org/}) is the main application used to process and visualise output from a \CNAME\ model. Install the \cname\ \R\ package which is part of the \CNAME\ download.

\CNAME\ has one \R\ package which is bundled with \CNAME\ application. This includes post processing function for extracting elements, in long form, from the Casal2 input for further analysis and plotting.  

There are three types of output that \CNAME\ can produce, depending on the type of analysis run. These outputs are: standard, MCMC, and tabular output.

The standard outputs are the reports that are produced in most \CNAME\ run modes, with the exception of \texttt{-s} and \texttt{-m}. The standard output can be split into two additional categories, a single parameter run (\texttt{casal2 -r}) or a multi-parameter run (\texttt{casal2 -r -i many\_pars.out}), or running in projection mode (\texttt{-f 1}). The standard output can be read into \R\ using the \texttt{extract.mpd()} function.

The second type of output is generated when doing an MCMC analysis (\texttt{casal2 -m}), which can generate two files, \texttt{mcmc\_objective.out} and \texttt{mcmc\_samples.out}. The MCMC outputs can be used to summarise convergence properties or chain behaviour, and can also be used to view marginal posteriors and quantify parameter uncertainty.

The third output type is the tabular outputs, either as space delimited (tabular) or as tab delimited (tabular-tsv). The tabular output can be generated after an MCMC analysis is done, to produce the marginal posteriors for various quantities, e.g., SSBs, model fits, etc. To get the posterior distributions for these derived quantities use the \texttt{-{}-tabular} or  \texttt{-{}-tabular-tsv} flag (e.g., \texttt{casal2 -r -i mcmc\_samples.out -{}-tabular > Tabular\_report.out}). This output can then be read into \R\ using the \texttt{extract.tabular()} function.

Note that the tabular-tsv will be much faster to read into \R\, \R\ nativity scans for tabs rather than spaces when reading a t4ext file.

\CNAME's reported output is written so that each \command{report} will start with a '*' and end with `*end'. This format can be used as the basis to construct functions that read \CNAME\ output to identify and read individual reports for post-processing.

The \CNAME\ \R\ \texttt{extract()} functions differ by how the expected output is structured and they each create a different \cname\ object. The \texttt{summary()} and \texttt{plot()} functions will generate different plots for the different \cname\ objects. Objects produced by the \texttt{extract()} function can be queried with \texttt{class(object)}.

The list of \cname\ \R\ functions include:

\begin{itemize}
	\item \texttt{extract.mpd()}, which parses the \CNAME\ default output into a list
	\item \texttt{extract.mcmc()}, which parses the \CNAME\ MCMC samples and objectives files into a list
	\item \texttt{extract.tabular()}, which parses the \CNAME\ tabular (space-separated or tab-separated) output into a list
	\item \texttt{extract.parameters()}, which parses the \CNAME\ free-parameter files into a list
	\item \texttt{extract.covariance.matrix()}, which reads the covariance matrix from an \texttt{mpd.out} estimation file
	\item \texttt{extract.correlation.matrix()}, which reads the correlation matrix from an \texttt{mpd.out} estimation file
	\item \texttt{extract.csl2.file()}, which reads a \CNAME\ \texttt{.csl2} (configuration) file into a list
	\item \texttt{generate.starting.pars()}, which reads in a file that contains the \command{estimate} blocks and generates \emph{N} starting values to test convergence
	\item \texttt{burn.in.tabular()}, which omits the first \emph{N} rows from a \texttt{casal2TAB} object
	\item \texttt{write.csl2.file()}, which writes a \CNAME\ \texttt{.csl2} (configuration) file to disk
	\item \texttt{ReadSimulatedData()}, which parses \CNAME\ output from a \texttt{casal2 -s} run
	\item \texttt{Method.TA1.8()}, which returns a weighting factor for age or length composition data. See \cite{francis2011data} for more detail
\end{itemize}

For the full list of functions available in the \CNAME\ \R\ package, including \texttt{get\_*()} helper functions for extracting specific objects from parsed output, refer to the package help system after loading the library (\texttt{library(Casal2)}; then \texttt{help(package = "Casal2")}).

The required and optional arguments for these functions can be queried after loading the \CNAME\ \R\ library with \texttt{library(Casal2)} and using the standard \R\ help syntax \texttt{?} (e.g., \texttt{?param.profile()}). Many of the help files have example code and data to demonstrate function syntax.

\paragraph*{Troubleshooting the \cname\ \R\ package}

If you get this error when using one of the \texttt{extract()} functions

\begin{lstlisting}
Read 1 item
Warning messages:
1: In scan(filename, what = "", sep = "\n", fileEncoding = fileEncoding) :
embedded nul(s) found in input
2: In extract.mpd(file = "results.txt", fileEncoding = "") :
File is empty, no reports found
\end{lstlisting}

You may be able to resolve this issue by using an alternative UTF format by specifying this format with the \subcommand{fileEncoding} parameter

\begin{lstlisting}
MyOutput <- extract.mpd(file = "Estimate.log", path = getwd(), fileEncoding = "UTF-16LE")
\end{lstlisting}
